% =============================================================================
% Part VII: What CausalFade Teaches Us (Chapters 26--29)
% =============================================================================
\chapter{Lessons in Scientific Research}\label{ch:lessons_research}

For students and fellow researchers embarking on empirical AI investigations, Project CausalFade offers eight concrete methodological lessons---none of them abstract platitudes, each paid for in real debugging hours across the repository's 320 commits:

\begin{enumerate}
  \item \textbf{Validate Your Estimators on Synthetic Truth Before Querying Real Models (\texttt{design\_validation}).} Running $3{,}600$ synthetic EM fits \emph{before} Cohort~001 cost minutes of CPU time and immediately revealed that the 3-parameter $C/A/B$ tree was observationally equivalent to symmetric logistic integration ($p=0.433$) and unstable near $C \to 1$. Pre-experiment simulation tells you what your statistical microscope can and cannot resolve.
  \item \textbf{Audit Raw Completions Before Trusting Aggregate Scores (Cohorts~001, 013, 024, 026).} In Cohort~001, \texttt{0/64} strict accuracy hid valid JSON inside Markdown fences. In Cohort~013, \texttt{0/64} hid valid JSON followed by trailing causal essays. In Cohort~024, \texttt{Llama Q4\_K\_M}'s control drop (\texttt{78/256}) was $448$ truncated prose explanations rather than wrong JSON choices. In Cohort~026, \texttt{Llama v2}'s \texttt{48/112} validity was a single trailing period (\texttt{"X=1."}). \textbf{Always inspect raw responses and separate format validity from semantic accuracy.}
  \item \textbf{Single-Item Binary Accuracy Is Blind to Constant Priors and Positional Policies (Cohorts~005, 024, 029).} Any binary benchmark with balanced labels grants $50\%$ accuracy and $0\pp$ conflict penalty to a constant \texttt{activate=1} policy (\texttt{995/1,024} in Cohort~029) or a \texttt{pick-first-option} policy (\texttt{256/256} in Cohort~024 \texttt{Mistral F16}). Always report directional class accuracies (\texttt{Gold=0} vs.\ \texttt{Gold=1}), counterbalanced option-flip consistency, and pairwise bundle fidelity ($F_{\mathrm{pair}} = I_0 I_1$).
  \item \textbf{Be Most Skeptical of Your Most Spectacular $p$-Value (Cohort~016).} Cohort~016 achieved $t(95)=21.10$, $p = 1.24\times 10^{-37}$ on a pre-registered 96-family replication---the dream statistic of many empirical papers. Yet cross-scale testing in Cohort~017 (\texttt{14B}: $0\pp$), cross-scaffold testing in Cohort~024 ($0\pp$), and retrospective magnitude auditing ($768/768$ solvable by picking the smaller printed gap) showed how fragile and shortcut-ridden that benchmark actually was. \textbf{Tiny $p$-values measure random sampling noise across worlds, NOT construct validity.}
  \item \textbf{Never Attribute Failure to a High-Level Construct Without a Matched Low-Level Control (Cohorts~002, 023, 028, 029).} Whenever an LLM fails a ``causal reasoning'' prompt, ask first whether it passes the exact same numerical operation stripped of causal framing. Cohort~023 (\texttt{58.6\%} harmful-direction control), Cohort~028 (\texttt{45.3\%} mental named-rule math vs.\ \texttt{100\%} with calculator), and Cohort~029 (\texttt{28.52\%} arithmetic pair fidelity) proved that baseline multi-number aggregation is a massive confound in single-pass LLM evaluation.
  \item \textbf{Pre-Declare Reference Competence Gates---and Obey Them When They Fail (Cohorts~004, 012, 018, 019, 026--028).} Stopping an experiment when the unquantized baseline fails (\texttt{stopped\_by\_gate}) feels painful in the moment, but it prevents polluting literature with meaningless comparisons between two broken systems.
  \item \textbf{Never Delete Failed Pilots, Interrupted Logs, or Buggy Analyzer Versions (Cohorts~008, 009--011, 017, 026).} Preserving \texttt{cohort\_008/run\_001\_f16} (\texttt{26/64}), \texttt{source\_order\_analysis\_v1\_undercounted\_counts.json}, Cohort~017's missing-NumPy trace, and \texttt{cohort\_026/invalidated\_pilot\_qwen\_v1} (\texttt{112} calls) with explicit README explanations created an unassailable audit trail.
  \item \textbf{Automate End-to-End Number Extraction From Raw Cryptographic Logs (\texttt{manuscript/scripts/extract.py}).} The \emph{only} numerical discrepancy found in the entire Cohort~029 release (\texttt{121/512, 416/512} vs.\ \texttt{14/512, 497/512} in \texttt{RESULTS.md}) occurred where prose was typed manually rather than generated from \texttt{events.jsonl}. Machine-verified extraction pipelines eliminate human transcription error.
\end{enumerate}

% =============================================================================
\chapter{Lessons in Causal AI}\label{ch:lessons_causal_ai}

Beyond general research methodology, CausalFade yields four sharp substantive insights into how current open-weight instruction-tuned transformers ($3\text{B}$--$14\text{B}$ parameters) interact with causal inference problems:

\section{1. Declarative Causal Knowledge $\neq$ Procedural Causal Execution}
Across multiple cohorts, models demonstrated near-perfect \emph{declarative} familiarity with textbook causal rules when queried in isolation: \texttt{Qwen2.5-7B F16} named the valid backdoor adjustment set \texttt{["Z"]} on $32/32$ ($100\%$) of Cohort~004 prompts and identified the randomized study as the governing causal source on $128/128$ ($100\%$) of Cohort~021--022 prompts, while \texttt{Llama-3.1-8B F16} answered direct causal graph ancestry questions at $31/32$ ($96.9\%$) in Cohort~027. Yet those exact same models failed completely to \emph{bind} that declarative rule to numerical selection and weighting in a single forward pass ($17/32$ actions in Cohort~004; $58/128$ guided actions in Cohort~021; $0/8$ graph switches in Cohort~027). Current instruction-tuned LLMs store verbal rules about causality without reliably compiling them into single-pass computational routing over tabular numbers.

\section{2. Tool Calculators Solve Estimation Arithmetic---Not Causal Identification}
A popular assumption in agentic AI is that coupling an LLM to a Python interpreter or calculator solves quantitative causal inference. Cohort~028 tested that assumption directly in a zero-leakage controlled ablation: an allowlisted deterministic Python calculator (\texttt{System C}) vaulted \texttt{Qwen2.5-7B F16} from $45.3\%$ to \textbf{$100.0\%$ ($128/128$)} on matched noncausal tasks where the rule name (\texttt{marginal} vs.\ \texttt{conditional}) was stated explicitly, yet yielded \textbf{$53.1\%$ ($68/128$, $+2.34\pp$)} on causal confounder-vs.-mediator tasks (\texttt{46.9\%} in \texttt{Llama-3.1-8B F16})! Externalizing calculation rescues \emph{statistical estimation} only if the agent can first solve \emph{structural causal identification}---mapping temporal/graphical structure onto the right adjustment functional \citep{pearl2009causality,sawarni2026crb,leban2026causalds}.

\section{3. Causal Framing Can Trigger Stronger Action Bias Than Neutral Framing}
One of Cohort~029's most surprising empirical discoveries was that \texttt{Qwen2.5-7B Q4\_K\_M} behaved \emph{more} degenerately under natural-prose causal advising (\texttt{995/1,024 = 97.17\%} \texttt{activate=1}; \texttt{3/256} target switches; \texttt{3/256} source-role switches) than under matched noncausal scorecard wording on the exact same numbers (\texttt{597/1,024 = 58.30\%} \texttt{activate=1}; \texttt{27/256} target switches; \texttt{67/256} source-role switches)! Why? Advisory prompts asking \emph{``whether switching a certain setting on would raise the success rate''} appear to activate strong RLHF/instruction-tuning helpfulness or affirmative-intervention priors (\texttt{activate=1}) whenever multi-table uncertainty exceeds single-pass capacity.

\section{4. Transportability and Pairwise Counterfactual Perturbations Are the Right Testbed}
Static observational question answering conflates memorized domain priors with causal reasoning, while two-study aggregate comparisons collapse into trivial single-pair inequalities once a study is chosen. Target-population transportability under subgroup heterogeneity (\texttt{CF-TP1})---requiring an agent to reverse its recommendation when target composition crosses $w_{\mathrm{crit}}$~\eqref{eq:w_crit} and retain its recommendation when it does not---provides an algebraic, domain-neutral litmus test for genuine interventional competence.

% =============================================================================
\chapter{The Unanswered Questions}\label{ch:unanswered}

Because Project CausalFade prioritized methodological rigor over overclaiming, its conclusion opens four crisp, directly testable research programs. For each open question, we specify the exact experimental design and controls required so future investigators do not repeat the pitfalls uncovered in Cohorts~001--028.

\section{1. Does Chain-of-Thought or Scratchpad Generation Unlock \texttt{CF-TP1} and Graph Switching?}
\begin{itemize}
  \item \textbf{The Open Question:} Cohort~029 (\texttt{CF-TP1}) and Cohorts~027--028 constrained models to emit bare tokens (\texttt{X=1}) or single-key JSON (\texttt{\{"activate":1\}}) in a single stateless pass without intermediate reasoning tokens. Since computing $w^* \mu_{\mathrm{rct}}(0,a) + (1-w^*)\mu_{\mathrm{rct}}(1,a)$ requires multi-step arithmetic and source filtering, how much of Cohort~029's failure disappears when explicit chain-of-thought scratchpads or reasoning-tuned test-time compute are permitted?
  \item \textbf{Proposed Controlled Design:} Evaluate \texttt{CF-TP1}'s exact frozen $4{,}096$-prompt fixture under a crossed $2\times 2$ generation regime: \emph{(no-scratchpad single JSON vs.\ structured scratchpad + final JSON)} $\times$ \emph{(causal transport vs.\ matched arithmetic scorecard)}, with generous \texttt{max\_tokens} ($\ge 1{,}024$) to prevent \texttt{Lian et al.} token-inflation truncation.
  \item \textbf{Diagnostic Criterion:} If scratchpad reasoning raises \emph{both} arithmetic and causal target-pair fidelity to $>85\%$, Cohort~029's collapse was a single-pass working-memory bottleneck; if scratchpads raise arithmetic to $>90\%$ while causal remains stalled (mirroring Cohort~028's calculator split!), the bottleneck is genuinely causal provenance/transport identification.
\end{itemize}

\section{2. Does \texttt{H01} (Selective Quantization Decay) Emerge Once Reference Models Pass \texttt{CF-TP1}?}
\begin{itemize}
  \item \textbf{The Open Question:} CausalFade could never test selective quantization decay of genuine causal transport because \texttt{3B}--\texttt{8B} single-pass models sat at the $25\%$ floor even at \texttt{F16}. What happens when a model/scaffold combination \emph{actually clears} $\ge 80\%$ switch-pair and invariant-pair fidelity at \texttt{F16}/\texttt{BF16} (e.g., \texttt{32B}--\texttt{70B} models or scratchpad-enabled \texttt{7B}--\texttt{14B} reasoning models) and is then compressed across \texttt{Q8\_0}, \texttt{Q6\_K}, \texttt{Q4\_K\_M}, \texttt{Q3\_K\_M}, and \texttt{Q2\_K} (both with and without \texttt{imatrix} calibration and QAT)?
  \item \textbf{Proposed Controlled Design:} Pre-register a reference gate requiring $\ge 80\%$ switch-pair fidelity and $\ge 80\%$ invariant-pair fidelity on \emph{both} causal and arithmetic \texttt{CF-TP1} splits at \texttt{BF16}/\texttt{F16}, then estimate the double difference $\Delta_{\mathrm{PTQ}}(p) = [F_{\mathrm{causal}}(p) - F_{\mathrm{causal}}(\text{F16})] - [F_{\mathrm{arith}}(p) - F_{\mathrm{arith}}(\text{F16})]$ across family clusters.
\end{itemize}

\section{3. Why Do Models Fail Confounder-vs.-Mediator Rule Selection With a Calculator (Cohort~028)?}
\begin{itemize}
  \item \textbf{The Open Question:} In Cohort~028, \texttt{Qwen 7B F16} and \texttt{Llama 8B F16} achieved $97\%$--$100\%$ direct graph-ancestry reading yet chose \texttt{marginal} vs.\ \texttt{conditional} weighting at $\sim 50\%$ chance even when both candidate definitions were written in the prompt. Is the failure caused by unfamiliar terminology (\texttt{marginal} vs.\ \texttt{conditional}), inability to trace backdoor vs.\ frontdoor/mediator causal logic, or confusion when both $P(W)$ and $P(W\mid X)$ appear simultaneously in the prompt?
  \item \textbf{Proposed Controlled Design:} Extend Cohort~028's \texttt{System C} calculator pipeline with a multi-stage Socratic decomposition: Stage~1 asks \emph{``Does X cause W, or does W cause X?''}; Stage~2 asks \emph{``If we intervene to set X=x, does the distribution of W stay at P(W) or change to P(W|X=x)?''}; Stage~3 dispatches the chosen distribution to the deterministic calculator.
\end{itemize}

\section{4. Mechanistic Localization Once Behavioral Competence Is Verified}
\begin{itemize}
  \item \textbf{The Open Question:} CausalFade rightly foreswore ``circuit severing'' claims (\texttt{D01}) because behavioral competence was absent. On a model that \emph{passes} \texttt{CF-TP1} ($>85\%$ switch and invariant fidelity), how do attention heads bind provenance tags (\texttt{randomized study}) to downstream stratum tables while suppressing confounded distractor tables?
  \item \textbf{Proposed Controlled Design:} Perform causal mediation / activation patching across \texttt{CF-TP1}'s minimal role-swap pairs (prompts differing \emph{only} by swapping the words \texttt{randomized} and \texttt{historical} between Source~A and Source~B on identical numerical tables), localizing the exact layers where provenance flips reverse the output logit difference.
\end{itemize}

% =============================================================================
\chapter{Epilogue}\label{ch:epilogue}

Twenty-nine days is a brief span on a calendar, yet in the life of a rigorously conducted research project it can encompass an entire scientific education. Project CausalFade set out on October 28, 2026 hoping to measure how weight compression erodes causal reasoning in local language models. Along the way, it encountered nearly every trap waiting for modern AI evaluators: Markdown fences masquerading as model incompetence, 2-bit token loops, psychometric models that looked principled on paper yet violated additivity on real outputs, numeric estimates that disagreed with a model's own actions on the very same line, a $p = 1.24 \times 10^{-37}$ prospective confirmation undermined by an unnoticed numerical gap invariant, $100\%$ first-option positional policies, and 4-bit verbosity truncations that struck control prompts while sparing causal ones.

What makes CausalFade worthy of study is not that those traps appeared---they appear silently in dozens of published LLM benchmarks every month---but that \textbf{every time the evidence contradicted an appealing narrative, the researchers listened to the evidence}. They did not delete failed cohorts. They did not hide interrupted runs or post-run script bugs. They did not promote an ungeneralizable $7\text{B}$ order swing as a universal law of compression, nor did they cherry-pick a confounded $6/8$ graph-switch run when counterbalancing dropped it to $0/8$. And when their culminating $4{,}096$-prompt confirmatory study (\texttt{CF-TP1}) yielded an inconclusive $-3.52\pp$ difference between causal transport and matched weighted arithmetic, they reported the null plainly---while showing through pairwise reversal and retention analysis exactly \emph{how} a language model can score $50\%$ on items while failing $98.8\%$ of required causal reversals.

Science advances at least as much by sharpening our instruments and exposing false positives as by announcing triumphs. In trading an unidentifiable story about quantization fade for a permanent, shortcut-audited diagnostic methodology, Project CausalFade fulfilled the highest obligation of empirical inquiry: leaving the field clearer, harder to fool, and better equipped to tell genuine causal reasoning apart from its many imitations.
